% Part: second-order-logic
% Chapter: metatheory
% Section: second-order-arithmetic

\documentclass[../../../include/open-logic-section]{subfiles}

\begin{document}

\olfileid{sol}{met}{spa}

\olsection{حساب مرتبهٔ دوم}

به یاد آورید که نظریهٔ حساب پئانو، یعنی~$\Th{PA}$، هشت اصل
$\Th{Q}$ را در بر می‌گیرد:
\begin{align*}
& \lforall[x][\eq/[x'][\Obj 0]]\\
& \lforall[x][\lforall[y][(\eq[x'][y'] \lif \eq[x][y])]]\\
& \lforall[x][(\eq[x][\Obj 0] \lor \lexists[y][\eq[x][y']])]\\
& \lforall[x][\eq[(x + \Obj 0)][x]]\\
& \lforall[x][\lforall[y][\eq[(x + y')][(x + y)']]]\\
& \lforall[x][\eq[(x \times \Obj 0)][\Obj 0]]\\
& \lforall[x][\lforall[y][\eq[(x \times y')][((x \times y) + x)]]]\\
& \lforall[x][\lforall[y][(x < y \liff \lexists[z][\eq[(z' + x)][y]])]]\\
\intertext{به‌علاوهٔ همهٔ جمله‌های به‌شکل زیر:}
& (!A(\Obj 0) \land \lforall[x][(!A(x) \lif !A(x'))]) \lif \lforall[x][!A(x)].
\end{align*}
مورد اخیر یک «شِما» است، یعنی الگویی که بی‌نهایت !!{sentence}s از زبان
حساب تولید می‌کند: یکی برای هر !!{formula} چون~$!A(x)$. این شِما را
\emph{شمای اصل استقرای} (مرتبهٔ اول) می‌نامیم.
در حساب پئانوی \emph{مرتبهٔ دوم}، یعنی~$\Th{PA^2}$، می‌توان استقرا را
با یک جملهٔ واحد بیان کرد. $\Th{PA^2}$ از هشت اصل نخست بالا به‌علاوهٔ
\emph{اصل استقرای} (مرتبهٔ دوم) تشکیل می‌شود:
\[
\lforall[X][((X(\Obj 0) \land \lforall[x][(X(x) \lif X(x'))]) \lif \lforall[x][X(x)])].
\]
این اصل می‌گوید اگر زیرمجموعه‌ای چون~$X$ از !!{domain}،
$\Assign{\Obj{0}}{M}$ را در بر گیرد و همراه با هر~$x \in \Domain{M}$،
$\Assign{\prime}{M}(x)$ را نیز در بر گیرد (یعنی «نسبت به جانشین بسته»
باشد)، آنگاه همه‌چیز را در !!{domain} در بر می‌گیرد (یعنی
  $X = \Domain{M})$.

اصل استقرا تضمین می‌کند هر !!{structure}ی که آن را ارضا می‌کند تنها
!!{element}sی از~$\Domain{M}$ را در بر گیرد که اصول وجودشان را ایجاب
می‌کنند؛ یعنی ارزش‌های~$\num{n}$ برای~$n \in \Nat$. مدل~$\Th{PA^2}$
هیچ عدد نااستانداردی ندارد.

\begin{thm}
\ollabel{thm:sol-pa-standard}
اگر~$\Sat{M}{\Th{PA^2}}$، آنگاه
$\Domain{M} = \Setabs{\Value{\num{n}}{M}}{n \in \Nat}$.
\end{thm}

\begin{proof}
بگذارید~$N = \Setabs{\Value{\num{n}}{M}}{n \in \Nat}$ و فرض کنید
$\Sat{M}{\Th{PA^2}}$. آشکار است که برای هر~$n \in \Nat$،
$\Value{\num{n}}{M} \in \Domain{M}$؛ پس~$N \subseteq \Domain{M}$.

اکنون شمول در جهت دیگر را در نظر بگیرید. تخصیص متغیری چون~$s$ را با
$s(X) = N$ در نظر بگیرید. بنا بر فرض،
\begin{align*}
\Sat{M}{\lforall[X][((X(\Obj 0) \land \lforall[x][(X(x)
      \lif X(x'))]) \lif \lforall[x][X(x)])]}, & \text{ پس}\\
\Sat{M}{(X(\Obj 0) \land \lforall[x][(X(x) \lif X(x'))]) \lif
  \lforall[x][X(x)]}[s]. &
\end{align*}
مقدم این شرطی را در نظر بگیرید. $\Value{\Obj{0}}{M} \in N$ و بنابراین
$\Sat{M}{X(\Obj{0})}[s]$. عطف دوم،
$\lforall[x][(X(x) \lif X(x'))]$، نیز ارضا می‌شود. زیرا فرض کنید
  $x \in N$. بنا بر تعریف~$N$، $x = \Value{\num{n}}{M}$ برای عددی
  چون~$n$. از این نتیجه می‌شود
$\Assign{\prime}{M}(x) = \Value{\num{n+1}}{M} \in N$. پس
$\Assign{\prime}{M}(x) \in N$.

داریم~$\Sat{M}{X(\Obj 0) \land \lforall[x][(X(x) \lif X(x'))]}[s]$.
در نتیجه~$\Sat{M}{\lforall[x][X(x)]}[s]$. اما این بدان معناست که برای
هر~$x \in \Domain{M}$ داریم~$x \in s(X) = N$. پس
$\Domain{M} \subseteq N$.
\end{proof}

\begin{cor}
\ollabel{cor:sol-pa-categorical}
هر دو مدل~$\Th{PA^2}$ هم‌ریخت‌اند.
\end{cor}

\begin{proof}
بنا بر~\olref{thm:sol-pa-standard}، دامنهٔ هر مدل~$\Th{PA^2}$ با
$\Value{\num{n}}{M}$ها پر می‌شود. هر چنین مدلی، مدل~$\Th{Q}$ نیز هست.
بنا بر~\olref[mod][mar][stm]{prop:thq-standard}، هر چنین مدلی استاندارد
است؛ یعنی با~$\Struct{N}$ هم‌ریخت است.
\end{proof}

در بالا~$\Th{PA^2}$ را نظریه‌ای تعریف کردیم که هشت اصل حسابی نخست و
اصل استقرای مرتبهٔ دوم را در بر دارد. در واقع، به‌لطف توان بیانی منطق
مرتبهٔ دوم، تنها \emph{دو اصل نخست} حسابی به‌همراه استقرا برای حساب
پئانوی مرتبهٔ دوم کافی‌اند.

\begin{prop}\ollabel{prop:sol-pa-definable}
بگذارید~$\Th{PA^{2\dagger}}$ نظریهٔ مرتبهٔ دومی باشد که دو اصل حسابی
نخست (اصول جانشین) و اصل استقرای مرتبهٔ دوم را در بر دارد. آنگاه
$\le$، $+$ و~$\times$ در~$\Th{PA^{2\dagger}}$ تعریف‌پذیرند.
\end{prop}

\begin{proof}
برای نشان‌دادن تعریف‌پذیری~$\le$ باید فرمولی چون~$!A_\le(x, y)$ بیابیم
که~$\Sat{N}{!A_\le(\num n, \num m)}$ اگر و تنها اگر~$n \le m$. فرمول
زیر را در نظر بگیرید:
\[
!B(x, Y) \ident Y(x) \land \lforall[y][(Y(y) \lif Y(y'))]
\]
آشکار است که~$!B(\num n, Y)$ به‌وسیلهٔ مجموعه‌ای چون
$Y \subseteq \Nat$ ارضا می‌شود اگر و تنها اگر
$\Setabs{m}{n \le m} \subseteq Y$؛ پس می‌توانیم بگیریم
$!A_\le(x, y) \ident \lforall[Y][(!B(x, Y) \lif Y(y))]$.

برای دیدن تعریف‌پذیری جمع، ملاحظه کنید که~$k+l=m$ اگر و تنها اگر تابعی
چون~$u$ وجود داشته باشد که~$u(0)=k$، $u(n')=u(n)'$ برای همهٔ~$n$،
  و~$m = u(l)$. می‌توانیم این هم‌ارزی را برای تعریف جمع
در~$\Th{PA^{2\dagger}}$ با فرمول زیر به کار بریم:
\[
!A_+(x,y,z) \ident \lexists[u][(u(\Obj 0)=x \land \lforall[w][u(w')=u(w)']
  \land u(y)=z)]
\]
باید روشن باشد که~$\Sat{N}{!A_+(\num k, \num l, \num m)}$ اگر و تنها
اگر~$k+l=m$.
\end{proof}

\begin{prob}
اثبات~\olref[sol][met][spa]{prop:sol-pa-definable} را کامل کنید.
\end{prob}

\end{document}
